\section[Notebooks 1968-15 and 1968-16: perfect subgroups, the elementary group and Whitehead's lemma]{Notebooks 1968-15 and 1968-16: perfect subgroups, the elementary group and Whitehead's lemma\protect\verifie{Notebook1968_15}}\label{sec:quillen-1968-15}

\emph{Context.} \carnet{1968}{15} and \carnet{1968}{16} contain the fair copy of the \q{plus} construction. In \S~1 of the first, \q{Killing a perfect subgroup of the fundamental group}, Proposition~1.2 kills a perfect subgroup $E$ of $\pi_1 X$; its proof (PDF, pp.~28--29) first replaces $E$ by the normal subgroup $N$ it generates, still perfect because $[N, N]$ is normal and contains $[E, E] = E$. In \S~2 of the second, \q{The K-theory associated to a ring} (p.~1), the construction is applied to the subgroup $E(A)$ of $GL(A)$ generated by the elementary matrices, which is perfect (Quillen refers to Bass); page~2 uses that $E(A)$ is normal in $GL(A)$, which rests on Whitehead's lemma. The ring $A$ is arbitrary there, \q{not necessarily commutative}.

\emph{Intuition.} The elementary matrix $e_{ij}(a) = 1 + a\,e_{ij}$ is a commutator as soon as a third index $k$ is available: $[e_{ik}(a), e_{kj}(b)] = e_{ij}(ab)$. For Whitehead's lemma, the block unipotent matrices $\left(\begin{smallmatrix} 1 & B \\ 0 & 1 \end{smallmatrix}\right)$ are products of elementary matrices, one for each coefficient of $B$, and three of them build the antidiagonal matrix $\left(\begin{smallmatrix} 0 & \theta \\ -\theta^{-1} & 0 \end{smallmatrix}\right)$.

\begin{encadre}
\begin{theoreme}[Quillen, notebooks 1968-15, p.~28, and 1968-16, pp.~1--2]\label{thm:quillen-1968-15}
Let $G$ be a group and $A$ a unital ring.
\begin{enumerate}
\item[a)] The normal subgroup generated by a perfect subgroup of $G$ is perfect.
\item[b)] If $i$, $k$, $j$ are pairwise distinct, $[e_{ik}(a), e_{kj}(b)] = e_{ij}(ab)$; in particular, the subgroup $E(n, A)$ of invertible matrices generated by the $e_{ij}(a)$ is perfect as soon as $n \ge 3$.
\item[c)] (Whitehead's lemma.) For $\theta$ invertible in $M_n(A)$, the matrix $\operatorname{diag}(\theta, \theta^{-1})$ is in $E(2n, A)$; for $\theta$, $\eta$ invertible, $\operatorname{diag}(\theta\eta\theta^{-1}\eta^{-1}, 1)$ is too.
\end{enumerate}
\end{theoreme}
\end{encadre}

\begin{proof}
a) Let $H$ be perfect and $N$ the normal subgroup it generates. The commutator subgroup $[N, N]$ is normal and contains $[H, H] = H$, hence contains $N$, so $[N, N] = N$.

b) Put $X = a\,e_{ik}$ and $Y = b\,e_{kj}$. Then $X^2 = Y^2 = YX = 0$, $XY = ab\,e_{ij}$, and $XY$ annihilates $X$ and $Y$ on the right; one checks that $(1 + X)(1 + Y)(1 - X)(1 - Y) = 1 + XY$. For $n \ge 3$, every generator $e_{ij}(a)$ is the commutator of $e_{ik}(a)$ and $e_{kj}(1)$, for a third index $k$.

c) $\left(\begin{smallmatrix} 1 & B \\ 0 & 1 \end{smallmatrix}\right)\left(\begin{smallmatrix} 1 & C \\ 0 & 1 \end{smallmatrix}\right) = \left(\begin{smallmatrix} 1 & B + C \\ 0 & 1 \end{smallmatrix}\right)$, and $\left(\begin{smallmatrix} 1 & a\,e_{ij} \\ 0 & 1 \end{smallmatrix}\right)$ is elementary; likewise for lower triangular matrices. A block computation gives
\[
\begin{pmatrix} 1 & \theta \\ 0 & 1 \end{pmatrix}\begin{pmatrix} 1 & 0 \\ -\theta^{-1} & 1 \end{pmatrix}\begin{pmatrix} 1 & \theta \\ 0 & 1 \end{pmatrix} = \begin{pmatrix} 0 & \theta \\ -\theta^{-1} & 0 \end{pmatrix},
\]
and the product of this matrix with the same one for $\theta = -1$ is $\operatorname{diag}(\theta, \theta^{-1})$. Finally
\[
\operatorname{diag}(\theta\eta\theta^{-1}\eta^{-1}, 1) = \operatorname{diag}(\theta, \theta^{-1})\operatorname{diag}(\eta, \eta^{-1})\operatorname{diag}((\eta\theta)^{-1}, \eta\theta).
\]
\end{proof}

\begin{sloppypar}
\noindent\emph{Verification in Lean.} a) is \texttt{normalClosure\_isPerfect}, b) is \texttt{commutator\_transvectionUnit} and \texttt{elementary\_isPerfect}, with the ring identity \texttt{commutator\_aux}, in \href{\leanblobmain/Notebook1968_15.lean}{\texttt{Notebook1968\_15.\allowbreak{}lean}}; c) is \texttt{exists\_mem\_elementary\_diag} and \texttt{exists\_mem\_elementary\_commutator}, with \texttt{upper\_mem\_elementary}, \texttt{lower\_mem\_elementary} and \texttt{upper\_lower\_upper}, in \href{\leanblobmain/Notebook1968_16.lean}{\texttt{Notebook1968\_16.\allowbreak{}lean}}. The matrices are the units of the ring \texttt{Matrix n n A}, for an arbitrary ring $A$.
\end{sloppypar}

\begin{sloppypar}\noindent Not formalized: the stable statements ($E(A)$ is normal in $GL(A)$, and $E(A)$ is the commutator subgroup $[GL(A), GL(A)]$), which require the inductive limit over $n$, and the \q{plus} construction itself.
\end{sloppypar}

\litterature{found in the literature: Milnor, \emph{Introduction to Algebraic K-Theory} (1971), \S~3 and Lemma~3.1; Rosenberg, \emph{Algebraic K-Theory and its Applications}, \S~2.1. mathlib has perfect groups (\texttt{Group.IsPerfect}) but none of these three statements. They are proposed for mathlib, and the edition does not present them as new.}
